\documentclass[a4paper,11pt]{ltjsarticle}
\usepackage{amsmath,amssymb,amsthm,mathtools,hyperref}
\newtheorem{theorem}{定理}[section]
\newtheorem{lemma}[theorem]{補題}
\newtheorem{proposition}[theorem]{命題}
\newtheorem{remark}[theorem]{注意}
\renewcommand{\proofname}{証明}
\newcommand{\qgr}{\operatorname{qgr}}
\newcommand{\gldim}{\operatorname{gldim}}
\newcommand{\thick}{\operatorname{thick}}
\newcommand{\Perf}{\operatorname{Perf}}
\newcommand{\Db}{D^{\mathrm b}}
\title{中心三次商の一般中心値と二つの中心チャート}
\author{}
\date{2026年11月21日・v4}
\begin{document}
\maketitle
\begin{abstract}
中心三次商について、二つの中心チャートによるperfectness判定と、
チャート代数が可換有限型中心部分環上有限である場合の一般中心値の正則性を扱う。
Rennemo--Segal--Van den Berghの非可換Bertini定理の有限中心性の仮定を保持し、
任意のPBW変形や特殊中心値へは適用しない。
一般点可換復元の正本は本repoの\texttt{papers/nc-cubic.tex}、
量子化の正本は別repo
\url{https://github.com/uedakazushi/cubic-quantization}である。
\end{abstract}

\section{二つの中心チャートによるperfectnessの判定}
\begin{proposition}[二つの中心チャート]\label{prop:patch}
$B$をconnected graded noetherian AS-regular代数、
$W\in Z(B)$を非零の正則斉次元、$A=B/(W)$とする。
$u,v\in Z(B)_1$は$B$上正則であり、
$A$上も中心的正則列をなすと仮定する。
$A[u^{-1}]_0$と$A[v^{-1}]_0$の大域次元は有限であり、
さらに境界について
\[
\Db(\qgr A/(u,v))
=\thick\{\pi(A/(u,v))(n)\mid n\in\mathbb Z\}
\]
を仮定する。このとき
\[
\Db(\qgr A)=\thick\{\pi A(n)\mid n\in\mathbb Z\}.
\]
\end{proposition}
\begin{proof}
$A$のcoherent moduleの有限生成graded lift $N$をとる。
十分高いsyzygyを取ってもline objectで生成される対象しか変わらないので、
$B$の有限大域次元と中心正則商のGorenstein性により、
$N$を有限生成自由加群から成るmatrix factorizationで表してよい。
$u$と$v$を局所化した二つのaffine chartでは
有限大域次元の仮定からfactorizationがcontractibleとなる。
contractionは有限個の係数・分母から成るので分母を払え、
ある$r,s\ge1$に対して
$u^r$と$v^s$の作用がprojective singularity categoryでは零になる。

二重Koszul cone $K(u^r,v^s;N)$のcohomologyは
$(u^r,v^s)$に殺される。
\[
(u,v)^{r+s-1}\subset (u^r,v^s)
\]
なので有限filtrationを持ち、各商が$A/(u,v)$-加群となる。
境界の生成性と中心正則列のKoszul分解より、
このconeはline objectのshiftから生成されるperfect対象である。
一方、projective singularity categoryでは$u^r,v^s$の作用は
零であるため、coneは$N$の四つのshift・twistの直和に分裂する。
従って$N$自身が直和因子としてperfectになる。
各coherent $N$に適用して主張を得る。
\end{proof}

\section{固定した周囲の一般中心三次商}
\begin{theorem}[有限中心チャートの一般二元三次係数]\label{thm:generic}
$B$をnoetherian connected graded AS-regular次元5の整域とする。
二つの次数1中心元$u,v$、三次中心元$G$を持ち、$B$が
$\mathbb C[u,v]$上自由であると仮定する。
$x,y,z$方向のfiltrationについて
\[
\operatorname{gr}_{x,y,z}B
\simeq S\otimes\mathbb C[u,v],
\]
$S$は整域、$G$の最高部分$\gamma\in S_3$は非零正則元とする。
係数空間$\mathbb A^4_{a,b,c,d}$の非空Zariski開集合$V$が存在し、
その各点について
\[
A_{a,b,c,d}
=B/(G+au^3+bu^2v+cuv^2+dv^3)
\]
が$\mathbb C[u,v]$上自由とする。
さらに
$\qgr B/(u,v,G)$が滑らかな楕円曲線のcoherent圏と同値であるとする。
さらに$D_u=B[u^{-1}]_0$、$D_v=B[v^{-1}]_0$がそれぞれ
可換有限型中心部分環$Z_u,Z_v$上有限生成加群であると仮定する。
$Z_u$は$v/u,G/u^3$を、$Z_v$は$u/v,G/v^3$を含むものとする。
このとき係数空間$\mathbb A^4_{a,b,c,d}$の非空Zariski開集合上で
\[
\gldim(\qgr A_{a,b,c,d})=3.
\]
\end{theorem}
\begin{proof}
$D_u=B[u^{-1}]_0$, $D_v=B[v^{-1}]_0$と置く。
$B$のAS-regular性と局所化から$D_u,D_v$は
noetherianで有限大域次元を持つ。
実際、$u$が中心次数1の単元で
$B[u^{-1}]\cong D_u[u,u^{-1}]$となり、
$D_u$-加群がLaurent拡大に誘導した後の
直和因子として回収できるからである。

$u$-chartに$t=v/u$を導入すると
\[
A_{a,b,c,d}[u^{-1}]_0
=D_u/(G/u^3+bt+ct^2+dt^3+a).
\]
ここで$a,b,c,d$は独立な普遍係数とする。
全係数族のchart代数は
\[
\mathcal D_u=D_u[a,b,c,d]/(G/u^3+bt+ct^2+dt^3+a)
             \simeq D_u[b,c,d]
\]
である。従って$\mathcal D_u$は有限大域次元を持ち、
可換有限型環$\mathcal Z_u=Z_u[b,c,d]$上有限である。
係数環$R=\mathbb C[a,b,c,d]$は
$a\mapsto-G/u^3-bt-ct^2-dt^3$によって中心へ作用する。
$Z_u$の有限個の$\mathbb C$-代数生成元は、この作用について
$\mathcal Z_u$を$R$上でも生成するので、$\mathcal Z_u$は有限型$R$-代数である。
最高active部分$\gamma$の非零性と整域性から、
この中心元の冪は$\mathbb C(b,c,d)$上一次独立であり、$R$の作用は忠実である。
Rennemo--Segal--Van den BerghのProposition~2.1を
$(R,\mathcal Z_u,\mathcal D_u)$へ適用すると、
非空開集合$U_u\subset\mathbb A^4$上の全chartが有限大域次元を持つ。
ここで重要なのは$\mathcal Z_u$上の加群有限性である。

同様に$v$-chartで$r=u/v$とすると
\[
A_{a,b,c,d}[v^{-1}]_0
=D_v/(G/v^3+ar^3+br^2+cr+d).
\]
$d$を消去した全族は$D_v[a,b,c]$である。
同じ命題を$Z_v[a,b,c]$上で用いれば非空開集合$U_v$を得る。
既に仮定した非空開集合$V$と$U_u\cap U_v$を交わらせる。
$\mathbb A^4$は既約なので、この共通部分も非空開集合である。
最高active部分の条件と$\mathbb C[u,v]$上の自由性により、
この範囲で$u,v$は$A_{a,b,c,d}$上の中心正則列である。
境界$C=B/(u,v,G)$の次数twistは有限生成graded moduleのpresentationにより
$\qgr C$を生成する。さらに$\qgr C\simeq\operatorname{coh}E$は
hereditaryであるため、twistの有限直和による二項presentation complexから
各coherent objectが直和因子として回収できる。
従って$\Db(\qgr C)=\thick\{\pi C(n)\}$であり、この境界の生成性と
命題\ref{prop:patch}を適用すれば
$\Db(\qgr A)=\thick\{\pi A(n)\}$。
最後に$B$のAS-regular性と中心三次正則商の
AS-Gorenstein次元4から、perfect tailsのSerre関手は
$T[3]$となる。ここで$T$はNakayama自己同型とGorenstein parameterに
対応する次数twistから成るexact自己同値である。
既に全coherent objectのperfectnessを示したので、任意の$M,N\in\qgr A$について
\[
 \operatorname{Ext}^p(M,N)
 \simeq\operatorname{Ext}^{3-p}(N,T(M))^*=0\qquad(p>3).
\]
また非零対象$M$について
$\operatorname{Ext}^3(M,T(M))\simeq\operatorname{End}(T(M))^*\ne0$
であるから、大域次元は正確に3である。
Gorenstein parameterの値を5と仮定する必要はない。
\end{proof}

\begin{remark}[適用限界と先行研究]
この定理は明記した有限中心性の下で、固定した周囲$B$に対する
非空Zariski開集合での十分条件を与える。
任意の特殊係数、一般の非PI Sklyanin因子、有限中心性を未確認の混合PBW変形には
そのまま適用できない。斉次特殊化がPIであっても、PBW変形の有限中心性は別途必要である。
主稿の定数型 $S(1,-1,-1)$ 混合族では、普遍六次中心元によって有限中心性を
別途証明した（主稿の定理\texttt{cmo:thm:finite-center}）。
その余有限Ku実現には射影的な中央対角と分離型への収縮を用いる。
一中心変数・二ブロックの64例外値について、この二中心変数の定理だけから
個別の滑らかさは得られない。

中心局所化・非可換generic smoothnessの外部入力として
J.~V.~Rennemo, E.~Segal, M.~Van den Bergh,
\emph{A non-commutative Bertini theorem},
\url{https://arxiv.org/abs/1705.01366}
のProposition~2.1とその証明（arXiv v2, pp.~3--6）を用いる。
同命題の仮定と、有限中心性が使われる箇所を一次資料で確認した。
Lemmas~3.2--3.3の局所化・分離体拡大だけでは、
一般の無限中心加群代数に対するsmoothnessの開性は結論できない。
この補助ノートの条件を、主稿の分離型の直接MF比較の仮定へ追加するものではない。
\end{remark}

\section{次に確認する入力}
混合表示へ適用するには、中心chartの可換有限型中心部分環とその上の有限生成性、
および境界のcoherent生成性を実際に示す必要がある。
これらを確認するまでは、中心係数のnormの根と特異係数の集合を同一視しない。
分離型の直接比較と混合型の現状は\texttt{00\_research\_plan.md}、
\texttt{notes/research\_notes.md}を参照する。
\end{document}
